\documentclass[12pt,a4paper]{article}

\usepackage[a4paper,margin=1.7cm]{geometry}
\usepackage[T2A]{fontenc}
\usepackage[utf8]{inputenc}
\usepackage[russian]{babel}
\usepackage{amsmath,amssymb,mathtools}
\usepackage{array,booktabs,longtable}
\usepackage{xcolor}
\usepackage{hyperref}
\usepackage{tikz}

\hypersetup{
    colorlinks=true,
    linkcolor=blue!55!black,
    urlcolor=blue!55!black
}

\setlength{\parindent}{0pt}
\setlength{\parskip}{0.55em}
\setlength{\tabcolsep}{4pt}
\renewcommand{\arraystretch}{1.25}

\newcommand{\tasktitle}[1]{%
  \phantomsection\label{task:#1}%
  {\LARGE\bfseries Задание №#1}\par
  \vspace{0.6em}
}

\begin{document}

\begin{titlepage}
\centering
\vspace*{0.22\textheight}

{\Huge\bfseries Ревью домашней работы\par}

\vspace{2.2cm}

{\Large Автор: Лёвин Артём Александрович\par}

\vspace{0.8cm}

{\large 21 сентября 2026 года\par}

\vfill

\begin{tikzpicture}
\draw[blue!35!black, line width=0.8pt]
  (0,0) .. controls (2,0.45) and (4,-0.45) .. (6,0)
        .. controls (8,0.45) and (10,-0.45) .. (12,0);
\end{tikzpicture}

\end{titlepage}

\newpage

\section*{Сводная таблица}

\small
\begin{longtable}{@{}p{1.2cm}p{2.4cm}p{5.0cm}p{3.3cm}p{3.3cm}@{}}
\toprule
\textbf{№ задания} & \textbf{Тема} & \textbf{Суть ошибки} & \textbf{Ваш ответ} & \textbf{Правильный ответ} \\
\midrule
\endfirsthead
\toprule
\textbf{№ задания} & \textbf{Тема} & \textbf{Суть ошибки} & \textbf{Ваш ответ} & \textbf{Правильный ответ} \\
\midrule
\endhead

\hyperref[task:6]{6} &
Неравенство треугольника &
Не учтена нижняя граница для третьей стороны: значения $1$, $2$ и $3$ не образуют треугольник со сторонами $5$ и $8$. &
$12$ &
$9$ \\

\addlinespace

\hyperref[task:14]{14} &
Геометрия: медиана и угол $30^\circ$ &
После верного равенства с числом $54$ в следующей строке число без основания заменено на $63$; решение не доведено до числового ответа. &
не указан &
$36$ \\

\addlinespace

\hyperref[task:16]{16} &
Среднее, медиана и мода &
Среднее и мода найдены верно, но для медианы выбраны не шестой и седьмой элементы упорядоченного ряда. &
среднее $3{,}5$; медиана $3{,}5$; мода $4$ &
среднее $3{,}5$; медиана $4$; мода $4$ \\

\bottomrule
\end{longtable}
\normalsize

\newpage
\vspace*{\fill}
\begin{center}
{\Huge\bfseries Разбор ошибок}
\end{center}
\vspace*{\fill}

\newpage
\tasktitle{6}

\textbf{Текст задания.}

Две стороны треугольника равны $5$ и $8$. Сколько целых положительных значений может принимать третья сторона?

\textbf{Ваше решение.}

В решении записаны три неравенства треугольника:
\[
a+b>c,\qquad b+c>a,\qquad a+c>b.
\]
После этого перечислены значения
\[
12,\ 11,\ 10,\ 9,\ 8,\ 7,\ 6,\ 5,\ 4,\ 3,\ 2,\ 1
\]
и сделан вывод, что возможны $12$ значений.

\textbf{Ваш ответ.}

$12$.

\textcolor{red}{\textbf{Ошибка:}} не применено строгое нижнее ограничение на третью сторону.

\textbf{Где именно возникает ошибка.}

Последний правильный шаг --- запись неравенств треугольника.

Первый неверный шаг --- включение в список допустимых значений чисел $1$, $2$ и $3$.

Обозначим третью сторону через $x$. Тогда для сторон $5$, $8$ и $x$ должны одновременно выполняться все три неравенства:
\[
5+8>x,\qquad 5+x>8,\qquad 8+x>5.
\]

Из первого получаем
\[
x<13.
\]

Из второго получаем
\[
x>3.
\]

Третье неравенство,
\[
8+x>5,
\]
выполняется для любого положительного $x$ и нового ограничения не добавляет.

\textbf{Почему это неверно.}

Неравенство треугольника требует, чтобы сумма длин любых двух сторон была \emph{строго} больше третьей стороны. Поэтому недостаточно проверить только верхнюю границу $x<13$.

Для значений $1$, $2$ и $3$ условие
\[
5+x>8
\]
нарушается:
\[
5+1=6<8,\qquad
5+2=7<8,\qquad
5+3=8.
\]
В последнем случае получается равенство, а не строгое неравенство. Треугольник с такими длинами сторон не существует.

\textcolor{blue!60!black}{\textbf{Ключевая идея:}} для третьей стороны треугольника с известными сторонами $a$ и $b$ нужно одновременно учитывать обе границы:
\[
|a-b|<x<a+b.
\]

\textbf{Исправление от последнего правильного шага.}

Из уже записанных Вами неравенств:
\[
5+8>x \quad\Rightarrow\quad x<13,
\]
\[
5+x>8 \quad\Rightarrow\quad x>3.
\]
Следовательно,
\[
3<x<13.
\]

Поскольку требуется найти только целые положительные значения, получаем:
\[
x=4,5,6,7,8,9,10,11,12.
\]

Таких значений девять.

\textbf{Полное правильное решение.}

Пусть третья сторона равна $x$.

По неравенству треугольника:
\[
5+8>x,
\]
откуда
\[
x<13.
\]

Также:
\[
5+x>8,
\]
откуда
\[
x>3.
\]

Неравенство
\[
8+x>5
\]
для положительного $x$ выполняется автоматически.

Значит,
\[
3<x<13.
\]

Целые положительные числа в этом промежутке:
\[
4,5,6,7,8,9,10,11,12.
\]

Их количество:
\[
9.
\]

\textbf{Проверка.}

Минимальное найденное значение $4$ подходит:
\[
5+4=9>8.
\]
Значение $3$ уже не подходит:
\[
5+3=8,
\]
а требуется строгое неравенство.

Максимальное найденное значение $12$ подходит:
\[
5+8=13>12.
\]
Значение $13$ уже не подходит, потому что
\[
5+8=13
\]
не больше $13$.

\textcolor{teal!60!black}{\textbf{Итог:}} третья сторона может принимать $9$ целых положительных значений.

\textbf{Задача на закрепление.}

Две стороны треугольника равны $7$ и $11$. Сколько целых положительных значений может принимать третья сторона?

\newpage
\tasktitle{14}

\textbf{Текст задания.}

В равнобедренном треугольнике $ABC$ выполнено
\[
AB=BC,\qquad \angle ABC=120^\circ.
\]
Медиана $BM$ проведена к основанию $AC$. На луче $BM$ за точкой $M$ выбрана точка $F$ так, что
\[
AF\perp AB.
\]
Известно, что
\[
FM=54.
\]
Найдите $AB$.

\textbf{Ваше решение.}

В решении правильно найдены углы основания:
\[
\angle BAC=\angle BCA=30^\circ.
\]
Также используется то, что медиана $BM$ в равнобедренном треугольнике является высотой и биссектрисой. Далее получены углы прямоугольного треугольника $ABF$ и записано, что при $AB=x$:
\[
BM=\frac{x}{2},\qquad BF=2x.
\]
Затем записано верное равенство
\[
BF=2x=\frac{x}{2}+54.
\]
Однако в следующих строках число $54$ заменено на $63$:
\[
\frac{x}{2}+63=2x,
\]
после чего числовой ответ не получен.

\textbf{Ваш ответ.}

не указан.

\textcolor{red}{\textbf{Ошибка:}} в алгебраическом преобразовании исходное число $54$ заменено на $63$.

\textbf{Где именно возникает ошибка.}

Последний правильный шаг:
\[
2x=\frac{x}{2}+54.
\]

Первый неверный шаг:
\[
\frac{x}{2}+63=2x.
\]

Число $63$ в условии и в предыдущем равенстве не появляется, поэтому такая замена не имеет математического основания.

\textbf{Почему это неверно.}

Длина $FM$ дана в условии и равна $54$. Поскольку точка $F$ лежит на луче $BM$ за точкой $M$, точки расположены на одной прямой в порядке
\[
B,\ M,\ F.
\]
Поэтому
\[
BF=BM+MF.
\]
Если $AB=x$, то из геометрии получается
\[
BM=\frac{x}{2},
\]
а значит,
\[
BF=\frac{x}{2}+54.
\]

Одновременно из прямоугольного треугольника $ABF$ получается
\[
BF=2x.
\]
Следовательно, именно равенство
\[
2x=\frac{x}{2}+54
\]
нужно решать дальше. Любая замена $54$ на другое число меняет условие задачи и приводит к другому результату.

\textcolor{blue!60!black}{\textbf{Ключевая идея:}} сначала связать $BM$ и $BF$ со стороной $AB$, а затем использовать точное равенство
\[
BF=BM+MF.
\]

\textbf{Исправление от последнего правильного шага.}

Продолжаем с верного равенства:
\[
2x=\frac{x}{2}+54.
\]

Умножим обе части на $2$:
\[
4x=x+108.
\]

Перенесём $x$ в левую часть:
\[
3x=108.
\]

Отсюда
\[
x=36.
\]

Значит,
\[
AB=36.
\]

\textbf{Полное правильное решение.}

Так как
\[
AB=BC,
\]
треугольник $ABC$ равнобедренный с основанием $AC$.

Угол при вершине $B$ равен
\[
120^\circ.
\]
Поэтому сумма двух равных углов при основании равна
\[
180^\circ-120^\circ=60^\circ.
\]
Следовательно,
\[
\angle BAC=\angle BCA=30^\circ.
\]

Медиана $BM$, проведённая к основанию равнобедренного треугольника, одновременно является высотой и биссектрисой. Поэтому
\[
BM\perp AC
\]
и
\[
\angle ABM=\frac{120^\circ}{2}=60^\circ.
\]

Рассмотрим прямоугольный треугольник $ABM$. В нём
\[
\angle AMB=90^\circ,\qquad \angle BAM=30^\circ.
\]
Катет $BM$ лежит напротив угла $30^\circ$, поэтому он равен половине гипотенузы $AB$:
\[
BM=\frac{AB}{2}.
\]

Теперь рассмотрим треугольник $ABF$. По условию
\[
AF\perp AB,
\]
значит,
\[
\angle BAF=90^\circ.
\]
Точки $B$, $M$, $F$ лежат на одной прямой, поэтому
\[
\angle ABF=\angle ABM=60^\circ.
\]
Тогда
\[
\angle AFB=30^\circ.
\]

В прямоугольном треугольнике $ABF$ сторона $AB$ лежит напротив угла $30^\circ$, поэтому она равна половине гипотенузы $BF$:
\[
AB=\frac{BF}{2},
\]
то есть
\[
BF=2AB.
\]

Положим
\[
AB=x.
\]
Тогда
\[
BM=\frac{x}{2},
\qquad
BF=2x.
\]

Так как $F$ лежит на луче $BM$ за точкой $M$, имеем
\[
BF=BM+MF.
\]
По условию
\[
MF=54,
\]
поэтому
\[
2x=\frac{x}{2}+54.
\]

Умножаем на $2$:
\[
4x=x+108.
\]
Следовательно,
\[
3x=108,
\]
откуда
\[
x=36.
\]

Значит,
\[
AB=36.
\]

\textbf{Проверка.}

Если
\[
AB=36,
\]
то
\[
BM=\frac{36}{2}=18.
\]
Тогда
\[
BF=BM+MF=18+54=72.
\]
С другой стороны, для треугольника $ABF$ должно выполняться
\[
BF=2AB=2\cdot36=72.
\]
Обе величины совпали.

\textcolor{teal!60!black}{\textbf{Итог:}} $AB=36$.

\textbf{Задача на закрепление.}

В равнобедренном треугольнике $ABC$ выполнено $AB=BC$ и $\angle ABC=120^\circ$. Медиана $BM$ проведена к основанию $AC$. На луче $BM$ за точкой $M$ выбрана точка $F$ так, что $AF\perp AB$. Известно, что $FM=72$. Найдите $AB$.

\newpage
\tasktitle{16}

\textbf{Текст задания.}

Оценки $2$, $3$, $4$, $5$ получили соответственно $2$, $3$, $6$, $1$ учащихся. Найдите среднюю оценку, медиану и моду.

\textbf{Ваше решение.}

Для среднего арифметического составлен ряд из всех оценок и получено:
\[
\frac{42}{12}=3{,}5.
\]

Для медианы записано:
\[
\frac{3+4}{2}=3{,}5.
\]

Для моды указано:
\[
4.
\]

\textbf{Ваш ответ.}

Среднее арифметическое --- $3{,}5$; медиана --- $3{,}5$; мода --- $4$.

\textcolor{red}{\textbf{Ошибка:}} неверно выбраны два центральных элемента при нахождении медианы.

\textbf{Где именно возникает ошибка.}

Последний правильный шаг --- упорядочивание оценок и вычисление среднего арифметического:
\[
\frac{42}{12}=3{,}5.
\]

Первый неверный шаг:
\[
\text{медиана}=\frac{3+4}{2}.
\]

Всего значений $12$, поэтому центральными являются шестой и седьмой элементы упорядоченного ряда. В данном ряду оба эти элемента равны $4$.

\textbf{Почему это неверно.}

Если количество чисел чётное, медиана равна среднему арифметическому \emph{двух центральных элементов} упорядоченного ряда.

Здесь ряд имеет вид:
\[
2,\ 2,\ 3,\ 3,\ 3,\ 4,\ 4,\ 4,\ 4,\ 4,\ 4,\ 5.
\]

Первые два места занимают двойки. Следующие три места занимают тройки, то есть после них заполнены места с первого по пятое. Поэтому шестое и седьмое места уже заняты четвёрками:
\[
\underbrace{2,\ 2,\ 3,\ 3,\ 3}_{5\text{ значений}},
\ \boxed{4},\ \boxed{4},
\ 4,\ 4,\ 4,\ 4,\ 5.
\]

Следовательно, нельзя брать числа $3$ и $4$: число $3$ находится на пятом месте, а не на шестом.

\textcolor{blue!60!black}{\textbf{Ключевая идея:}} при чётном количестве наблюдений сначала точно определить номера двух центральных позиций, а уже затем брать значения на этих позициях.

\textbf{Исправление от последнего правильного шага.}

Всего учащихся:
\[
2+3+6+1=12.
\]

Значит, для медианы нужны шестое и седьмое значения.

Упорядоченный ряд:
\[
2,\ 2,\ 3,\ 3,\ 3,\ 4,\ 4,\ 4,\ 4,\ 4,\ 4,\ 5.
\]

Шестое значение равно $4$, и седьмое значение тоже равно $4$. Поэтому:
\[
\text{медиана}=\frac{4+4}{2}=4.
\]

Среднее арифметическое $3{,}5$ и мода $4$ в Вашем решении найдены верно.

\textbf{Полное правильное решение.}

Общее количество оценок:
\[
2+3+6+1=12.
\]

Сумма всех оценок:
\[
2\cdot2+3\cdot3+4\cdot6+5\cdot1
=
4+9+24+5
=
42.
\]

Среднее арифметическое:
\[
\frac{42}{12}=3{,}5.
\]

Для медианы выпишем все значения по возрастанию:
\[
2,\ 2,\ 3,\ 3,\ 3,\ 4,\ 4,\ 4,\ 4,\ 4,\ 4,\ 5.
\]

Так как чисел $12$, медиана равна среднему арифметическому шестого и седьмого значений:
\[
\frac{4+4}{2}=4.
\]

Чаще всего встречается оценка $4$: она встречается шесть раз. Поэтому мода равна
\[
4.
\]

\textbf{Проверка.}

Количество выписанных значений действительно равно $12$.

Сумма значений в упорядоченном ряду равна $42$, поэтому среднее $3{,}5$ согласуется с исходными частотами.

Шестое и седьмое значения равны $4$, поэтому медиана равна $4$.

Наибольшая частота равна $6$ и относится к оценке $4$, поэтому мода также равна $4$.

\textcolor{teal!60!black}{\textbf{Итог:}} среднее арифметическое равно $3{,}5$, медиана равна $4$, мода равна $4$.

\textbf{Задача на закрепление.}

Оценки $2$, $3$, $4$, $5$ получили соответственно $1$, $4$, $5$, $2$ учащихся. Найдите среднюю оценку, медиану и моду.

\vfill
\begin{center}
\small Лёвин Артём Александрович эксклюзивно для Екатерины Скелли
\end{center}

\end{document}
